\chapter{導入}
\label{ch:introduction}

本資料は，Wasserstein 距離の定義から出発し，次の2つの結果を目標とする．
主たる文献は Villani
\emph{Optimal Transport: Old and New}
（Grundlehren der mathematischen Wissenschaften 338, Springer, 2009）であり，
定義・記法・証明はこれに従う．
Gaussian 測度の閉形式（目標2）は Givens--Shortt の論文
\emph{A class of Wasserstein metrics for probability distributions}
（Michigan Math.\ J.\ \textbf{31} (1984), 231--240）を原論文とする．
日本語の解説としては
桑江ほか『最適輸送理論とリッチ曲率』
（Encounter with Mathematics 第63回，中央大学，2015）の §1.4 を参照し，
測度分解定理・接着補題などの用語はこれに合わせた．

\begin{itemize}
  \item \textbf{目標1（距離性；Villani (2009) Chapter 6）}：
    Polish 空間 $\mathcal{X}$ 上で
    \[
      1\leq p<\infty
      \quad\Longrightarrow\quad
      \Wass_p\text{ は }\Pp{p}(\mathcal{X})\text{ 上の距離である．}
    \]
  \item \textbf{目標2（Gaussian 測度の閉形式；Givens--Shortt (1984)
    Proposition 7）}：
    $\R^n$ 上の Gaussian 測度に対して
    \[
      \Wass_2\bigl(\Normal(m_1,\Sigma_1),\,\Normal(m_2,\Sigma_2)\bigr)^2
      =
      \norm{m_1-m_2}^2
      +\tr\Bigl(\Sigma_1+\Sigma_2
        -2\bigl(\Sigma_1^{1/2}\Sigma_2\Sigma_1^{1/2}\bigr)^{1/2}\Bigr).
    \]
    右辺は平均のずれと共分散のずれに分離しており，
    Wasserstein 距離が閉形式で計算できる数少ない例である．
\end{itemize}

\section*{本資料の構成}

\begin{itemize}
  \item 第\ref{ch:wasserstein-metrics}章：
    $\Pp{p}(\mathcal{X})$，coupling，$\Wass_p$ を定義し，
    最適 coupling の存在定理（Villani (2009) Theorem 4.1）を経て
    目標1を証明する．
  \item 第\ref{ch:gaussian}章：$\R^n$ 上の Gaussian 測度を像測度として定義し，
    目標2を証明する（特異な共分散行列も含めた一般の場合を扱う）．
  \item 付録A：
    距離空間と測度論の準備．
    本編で用いる定義と補題をすべて揃え，
    証明なしで認める標準定理（単調収束定理・一致定理・分解定理など）は
    ステートメントを明示して文献を挙げる．
  \item 付録B：線形代数の準備（半正定値行列・平方根・特異値分解など）．
\end{itemize}

まず本編を順に読み，参照された定義・補題を必要に応じて
付録で確認する読み方を想定している．

\section*{主な記号}

\begin{itemize}
  \item $\Prob(\mathcal{X})$：$\mathcal{X}$ 上の Borel 確率測度全体．
  \item $\Pp{p}(\mathcal{X})$：有限 $p$ 次モーメントをもつ確率測度全体
    （定義~\ref{def:wp-p-space}）．
  \item $\Couplings(\mu,\nu)$：$\mu,\nu$ の coupling 全体
    （定義~\ref{def:wp-coupling}）．
  \item $\Wass_p(\mu,\nu)$：Wasserstein 距離（定義~\ref{def:wp-wasserstein}）．
  \item $\mathbf{1}_A$：指示関数（定義~\ref{def:meas-indicator}）．
  \item $\norm{f}_{L^p(\mu)}$：$L^p$ ノルム（定義~\ref{def:meas-lp-norm}）．
  \item $\mu\circ T^{-1}$：写像 $T$ による $\mu$ の像測度（定義~\ref{def:meas-pushforward}）．
  \item $\mathrm{pr}_1,\ \mathrm{pr}_{12}$ など：座標への射影（projection）．
    添字は残す成分の番号（例：$\mathrm{pr}_{12}(x,y,z)=(x,y)$）．
  \item $\mu\otimes\nu$：直積測度（定理~\ref{thm:meas-product}）．
  \item $\Normal(m,\Sigma)$：平均 $m$，共分散 $\Sigma$ の Gaussian 測度
    （第\ref{ch:gaussian}章）．
  \item $\tr A$，$\Sigma^{1/2}$，$A\succeq0$：
    トレース・半正定値行列の平方根・半正定値性（付録B）．
\end{itemize}

